
\chapter{Analyse des besoins et étude technique}

Dans ce chapitre, nous allons étudier l'environnement du recrutement et de la gestion documentaire au sein de l'entreprise afin de comprendre les différentes contraintes liées au traitement des CV. Ensuite, nous procéderons à l'analyse approfondie des besoins, en vue de les synthétiser et de proposer une architecture logicielle adaptée et efficace. Nous réaliserons par la suite une étude technique qui fournira une vue d'ensemble des technologies, des processus, et des outils évalués pour concevoir, développer et mettre en œuvre le projet AI-BioProfile. L'objectif principal de cette étude est de déterminer la faisabilité technique du projet, de justifier nos choix architecturaux et d'identifier les solutions les plus appropriées pour répondre aux exigences strictes de performance et de sécurité des données.

\section{Compréhension et recueil des informations sur les besoins}

\subsection{Compréhension de la nature des besoins}
Dans le cadre de mon projet de fin d'année, j'ai eu la responsabilité de conceptualiser et de réaliser une application complète d'intelligence artificielle pour automatiser la génération des dossiers de compétences (BioProfiles). Le processus de recrutement actuel implique la réception quotidienne de dizaines de CV sous des formats hétérogènes. La conversion manuelle de ces documents en présentations commerciales standardisées est une tâche à faible valeur ajoutée, source d'erreurs et de retards dans la proposition des candidats aux clients.

La solution à concevoir doit couvrir différents aspects du traitement de la donnée non structurée, en tenant compte des contraintes opérationnelles :
\begin{itemize}
    \item L'application doit être capable de lire et d'interpréter des formats de fichiers variés (PDF texte, PDF scannés, documents Word DOCX).
    \item L'extraction des données doit être sémantique (comprendre les expériences et compétences) et non basée sur de simples mots-clés.
    \item Le système doit garantir la confidentialité absolue des données des candidats, excluant de facto l'utilisation d'API d'intelligence artificielle hébergées sur le Cloud public.
    \item La restitution doit être immédiate et directement exploitable sous la forme d'un fichier PowerPoint respectant la charte graphique de l'entreprise.
\end{itemize}

\subsection{Recueil des informations sur les besoins}
Les réunions avec les différentes parties prenantes (Ressources Humaines, Ingénieurs d'Affaires) ont permis d'identifier et de segmenter les besoins en plusieurs modules architecturaux distincts :

\textbf{Module 1 : Pipeline d'ingestion et de prétraitement documentaire}
L'objectif est de préparer la donnée brute pour qu'elle soit lisible par l'intelligence artificielle. Cela nécessite de :
\begin{itemize}
    \item Concevoir un système d'extraction de texte natif et d'analyse d'images (Computer Vision) pour récupérer la photo de profil.
    \item Intégrer un système de Reconnaissance Optique de Caractères (OCR) robuste pour traiter les CV scannés.
\end{itemize}

\textbf{Module 2 : Pipeline d'inférence et d'Intelligence Artificielle}
L'objectif est d'analyser le texte extrait pour en dégager une structure de données claire. Cela implique de :
\begin{itemize}
    \item Déployer un modèle de langage (LLM) en local sur les serveurs de l'entreprise.
    \item Concevoir des instructions (Prompts) strictes pour forcer le modèle à renvoyer un format JSON précis (Nom, Titre, Compétences, Expériences).
\end{itemize}

\textbf{Module 3 : Pipeline de post-traitement et génération}
L'objectif est de convertir la donnée structurée en un livrable commercial. Cela nécessite de :
\begin{itemize}
    \item Développer un script capable de manipuler le format Office Open XML pour injecter du texte et des images dans un modèle PowerPoint existant.
\end{itemize}

\textbf{Module 4 : Interface Utilisateur et API}
L'objectif est de rendre le service accessible aux utilisateurs non techniques. Cela implique de :
\begin{itemize}
    \item Développer une API asynchrone capable de gérer plusieurs requêtes simultanées.
    \item Créer une interface web (Front-end) intuitive permettant le chargement (upload) des CV et la prévisualisation des données avant la génération finale.
\end{itemize}


\section{Identification des besoins}

Suite au recueil des informations, les exigences du projet ont été classifiées en deux catégories : les besoins fonctionnels (ce que le système doit faire) et les besoins non fonctionnels (comment le système doit se comporter).

\subsection{Besoins fonctionnels}
\begin{itemize}
    \item \textbf{Gestion de l'import de fichiers :} Le système doit permettre le téléchargement de fichiers aux formats PDF, DOCX et TXT.
    \item \textbf{Extraction multimodale :} Le système doit extraire non seulement le texte, mais aussi isoler et recadrer automatiquement la photo de profil du candidat présente sur le document.
    \item \textbf{Validation humaine dans la boucle (Human-in-the-loop) :} Avant de générer le fichier PowerPoint final, l'utilisateur doit pouvoir visualiser les données extraites par l'IA sur l'interface web, et corriger manuellement d'éventuelles erreurs ou omissions.
    \item \textbf{Génération du livrable :} Le système doit permettre le téléchargement direct du BioProfile au format .pptx, correctement formaté.
\end{itemize}

\subsection{Besoins non fonctionnels}
\begin{itemize}
    \item \textbf{Sécurité et Confidentialité (Privacy-by-design) :} C'est le besoin le plus critique. Les CV contiennent des données sensibles (RGPD). Aucune donnée textuelle ou image ne doit transiter par des serveurs externes. Les algorithmes d'IA doivent tourner en local (On-Premise).
    \item \textbf{Temps de réponse (Performance) :} L'ensemble du processus (du chargement du CV à la génération du PowerPoint) doit s'effectuer dans un délai raisonnable (idéalement inférieur à une minute) pour maintenir une expérience utilisateur fluide, malgré la lourdeur des traitements d'IA locale.
    \item \textbf{Résilience et fiabilité :} Le système doit être capable de détecter si un PDF est illisible (ex: scan de très mauvaise qualité) et renvoyer des messages d'erreur clairs à l'utilisateur sans faire crasher l'application.
\end{itemize}


\section{Étude Technique}

Afin de construire une architecture robuste répondant à ces besoins complexes, nous avons réalisé un benchmark des technologies disponibles sur le marché pour chaque brique de notre application. Les critères d'évaluation incluent la performance, la facilité d'intégration, le coût (licence), et surtout la capacité à fonctionner hors-ligne (localement).

\subsection{Benchmark et analyse des outils}

\textbf{1. Comparaison des outils d'extraction PDF}
Pour extraire le texte et les images des fichiers PDF, nous avons évalué :
\begin{itemize}
    \item \textit{PyMuPDF (fitz) :} Extrêmement rapide, écrit en C, il gère parfaitement les coordonnées spatiales du texte et l'extraction native des images.
    \item \textit{PyPDF2 :} Populaire et 100\% Python, mais nettement plus lent et souvent mis en échec par des formats de PDF récents.
    \item \textit{PDFMiner :} Très précis sur l'extraction de texte complexe, mais extrêmement lourd et lent, inadapté à un traitement quasi-temps réel.
\end{itemize}
\textit{Choix : PyMuPDF a été sélectionné pour ses performances supérieures et sa flexibilité dans l'extraction des images.}

\textbf{2. Comparaison des moteurs d'OCR (Reconnaissance Optique)}
\begin{itemize}
    \item \textit{Tesseract OCR :} Solution open-source soutenue par Google, très mature, supportant de nombreuses langues, fonctionnant parfaitement en local.
    \item \textit{EasyOCR :} Basé sur le Deep Learning (PyTorch), il est très précis mais nécessite une grande puissance de calcul (GPU) pour être rapide, ce qui l'alourdit.
    \item \textit{Google Cloud Vision API :} La meilleure précision du marché, mais éliminée d'office en raison de la contrainte de confidentialité des données (solution Cloud).
\end{itemize}
\textit{Choix : Tesseract OCR a été retenu pour son excellent compromis entre vitesse d'exécution sur CPU, gratuité et respect de la confidentialité.}

\textbf{3. Comparaison des moteurs d'Inférence IA (LLM Local)}
\begin{itemize}
    \item \textit{Ollama (avec modèle Mistral) :} Un outil open-source permettant d'exécuter facilement des grands modèles de langage en local. Le modèle Mistral offre des performances comparables à GPT-3.5 tout en étant assez léger pour tourner sur des machines standards.
    \item \textit{OpenAI API (ChatGPT) :} Très facile à implémenter et ultra-performant pour structurer du JSON, mais enfreint la règle de souveraineté des données exigée par le projet.
    \item \textit{Llama.cpp :} Très optimisé, mais plus complexe à mettre en place en tant qu'API de service (backend) par rapport à la simplicité d'Ollama.
\end{itemize}
\textit{Choix : Ollama couplé au modèle Mistral a été choisi car il permet une inférence 100\% locale, garantissant la sécurité des données tout en offrant une précision sémantique élevée.}

\textbf{4. Comparaison des Frameworks Backend Web (Python)}
\begin{itemize}
    \item \textit{FastAPI :} Framework moderne, nativement asynchrone, ultra-rapide. Il s'intègre parfaitement avec \textit{Pydantic} pour la validation stricte des données (JSON), ce qui est idéal pour communiquer avec un LLM.
    \item \textit{Flask :} Très simple et léger, mais la gestion asynchrone (nécessaire pour ne pas bloquer le serveur pendant que l'IA réfléchit) est moins native et moins performante que FastAPI.
    \item \textit{Django :} Framework extrêmement complet (Batteries-included), mais beaucoup trop lourd pour notre besoin qui se concentre sur la création d'une API de traitement de données.
\end{itemize}
\textit{Choix : FastAPI s'est imposé comme une évidence pour ses performances asynchrones et sa synergie avec Pydantic.}


\subsection{Définitions des Outils Choisis dans l'Architecture}

\textbf{FastAPI et Pydantic} \\
FastAPI est un framework web moderne et rapide pour la création d'API avec Python, basé sur les annotations de type standard de Python. Dans notre architecture, il agit comme le chef d'orchestre. Il réceptionne les requêtes des utilisateurs contenant les CV, distribue les tâches d'extraction aux différents modules, et renvoie le fichier PowerPoint final. Couplé à la bibliothèque \textit{Pydantic}, il nous permet de définir des schémas de données stricts (par exemple, obliger le LLM à retourner une liste d'expériences contenant spécifiquement une date, un titre et une description). Si le modèle d'IA tente de renvoyer une donnée mal formatée, Pydantic la rejette et force la validation.

\textbf{Ollama et le modèle Mistral} \\
Ollama est un outil open-source qui simplifie considérablement l'exécution locale de grands modèles de langage (LLMs). Il gère l'allocation mémoire et fournit une API REST locale. Dans notre pipeline, Ollama exécute le modèle franco-européen \textit{Mistral}. Ce modèle est sollicité pour lire le texte brut et désordonné du CV, en comprendre la sémantique, et structurer les données. Cette approche garantit que la "réflexion" de l'intelligence artificielle se fait à huis clos, sur la machine hôte, assurant une conformité totale en matière de protection des données (RGPD).

\textbf{PyMuPDF, Tesseract et OpenCV (Extraction et Vision)} \\
Ce triptyque forme notre moteur d'ingestion documentaire :
\begin{itemize}
    \item \textbf{PyMuPDF} parcourt le fichier PDF pour en extraire le texte natif très rapidement.
    \item Si le texte natif est insuffisant (cas d'un CV scanné), \textbf{Tesseract OCR} prend le relais pour "lire" les pixels de l'image et reconstituer le texte de secours.
    \item \textbf{OpenCV} (via les classificateurs de Haar) intervient en dernier recours pour analyser la géométrie de la page et détecter la présence d'un visage humain. Une fois détecté, le visage est recadré et sauvegardé temporairement pour devenir la photo de profil du document final.
\end{itemize}

\textbf{Python-pptx} \\
C'est une bibliothèque Python dédiée à la création et à la mise à jour de fichiers PowerPoint (.pptx). Dans notre architecture, elle représente la dernière étape du traitement. Le script charge un \textit{Template} (modèle vierge de l'entreprise), accède aux différentes diapositives (Slides) de manière programmatique (via l'XML sous-jacent), et remplace les balises fictives par les véritables données JSON validées (nom, compétences) ainsi que par la photo isolée précédemment.


\section{Conclusion}

L'analyse des besoins et l'étude technique réalisées dans ce chapitre ont permis de cerner avec précision les exigences complexes liées à l'automatisation de la gestion des CV. Nous avons identifié des contraintes fortes, particulièrement en matière de sécurité et de confidentialité, excluant le recours à la facilité des services Cloud. Grâce à un benchmark rigoureux, nous avons sélectionné une stack technologique performante et innovante : FastAPI pour l'orchestration asynchrone, PyMuPDF et Tesseract pour l'ingestion multimodale, Ollama/Mistral pour l'analyse intelligente locale, et python-pptx pour la génération du livrable. 

La faisabilité technique du projet AI-BioProfile est ainsi confirmée. Cette combinaison d'outils open-source garantit non seulement la protection des données, mais offre également les performances nécessaires pour transformer de manière transparente un document brut en un profil commercial qualitatif. En se basant sur ces fondations technologiques solides, le prochain chapitre se concentrera sur la phase de réalisation concrète, en détaillant l'architecture implémentée, les étapes de développement, et les résultats obtenus lors des phases de test.
```